This chapter has described the mechanism that produces the prices
which the rest of the course takes as given.

Three points are worth carrying forward.

First,
there is no such thing as \emph{the} price.
There is a best bid, a best ask, a mid-price, a last traded price,
and they differ.
Which one is appropriate depends on what is being computed,
and a great deal of sloppy quantitative work begins with the wrong choice.

Second,
liquidity is finite and it is consumed.
The size that can be traded at the best quote is a small fraction of a typical
institutional order,
so the cost of a trade is a function of its size,
not a constant per unit.

Third,
the book is a queueing system,
and the tools that describe it are counting processes and their intensities.
This is the natural point of contact between market microstructure and
the stochastic-process material of the course.

The accompanying notebook reconstructs a limit order book from a message file,
computes the quantities defined in Section \ref{sec.theBook},
and reproduces the linear relation between order flow imbalance and price change
discussed in Section \ref{sec.orderFlow}.
The exercise is a good test of whether the definitions above have been understood:
they are all easy to state and easy to get subtly wrong in code.

For a broad and readable treatment of the empirical facts, see \cite{Bou18tra};
for the classical theory of informed trading, see \cite{Kyl85con} and \cite{GM85bid};
and for a survey of limit order book models, see \cite{GPFA13lim}.
